\section{Three-Layer Corruption}
\label{sec:three-layer}

This section establishes the central empirical observation: unlearning side-effects decay across three semantic layers, but $\Lthree$ never reaches zero, and the layer profile depends substantially on which $\Df$ was chosen.

\subsection{Headline Numbers}
\label{sec:headline}

Table~\ref{tab:headline} reports geometric means of the per-sample PPL ratio $r$ across the $1.5$M evaluations.

\begin{table}[t]
\centering
\small
\setlength{\tabcolsep}{4pt}
\begin{tabular}{@{}lccc@{}}
    \toprule
                                                   & $\Lone$ & $\Ltwo$ & $\Lthree$ \\
    \midrule
    Samples                                        & $5{,}000$ & $10{,}000$ & $495{,}000$ \\
    Geo-mean $r$                                   & \textbf{$1.96\times$} & \textbf{$1.32\times$} & \textbf{$1.19\times$} \\
    \% samples with $r{>}1.1$                      & $100.0$ & $90.8$ & $73.8$ \\
    \% samples with $r{>}2.0$                      & $42.0$  & $2.4$  & $0.8$ \\
    \bottomrule
\end{tabular}
\caption{Three-layer corruption headline ($n{=}100$ forget sets, $1.5$M per-sample ratios). Decay is strict ($\Lone{>}\Ltwo{>}\Lthree$), but $\Lthree{>}1$ on $73.8\%$ of cross-topic samples.}
\label{tab:headline}
\end{table}

\begin{figure}[t]
\centering
\includegraphics[width=0.95\columnwidth]{fig_three_layer_decay.pdf}
\caption{\textbf{Three-layer corruption decay.} Left: geo-mean PPL ratio $r$ falls from $\Lone{=}1.96\times$ to $\Lthree{=}1.19\times$. Right: fraction of samples with $r{>}1.1$ ($100\%/90.8\%/73.8\%$). $\Lthree$ is non-zero in both senses.}
\label{fig:three-layer-decay}
\end{figure}

Two observations stand out.

\textbf{(C1) Strict monotone decay.} $\Lone{>}\Ltwo{>}\Lthree$ holds globally and within each cluster (Sec.~\ref{sec:invariants}); the layer ordering is not an artefact of aggregation.

\textbf{(C2) $\Lthree$ does not vanish.} Cross-cluster samples that are by construction unrelated to $\Df$ are corrupted on $73.8\%$ of measurements, with geo-mean $1.19\times$. Spillover is the rule, not a tail event.

\subsection{Per-Forget-Set Variance}
\label{sec:per-forget}

The same algorithm and hyperparameters applied to different forget sets produce noticeably different damage profiles. Table~\ref{tab:perforget} contrasts the mildest and worst forget set among $100$.

\begin{table}[t]
\centering
\small
\setlength{\tabcolsep}{4pt}
\begin{tabular}{@{}lccc@{}}
    \toprule
                                                   & $\Lone$ & $\Ltwo$ & $\Lthree$ \\
    \midrule
    Mildest triplet (\texttt{029})                 & $1.54$ & $1.15$ & $1.12$ \\
    Worst triplet (\texttt{073}, ``storm'')        & $\mathbf{2.44}$ & $\mathbf{1.73}$ & $\mathbf{1.26}$ \\
    \midrule
    Spread (max/min, $n{=}100$)                    & $\mathbf{1.59\times}$ & $1.50\times$ & $1.13\times$ \\
    \bottomrule
\end{tabular}
\caption{Forget-set-induced variation. $\Lone$ varies by $1.59\times$ between the mildest and worst $\Df$ under fixed algorithm/hyperparameters.}
\label{tab:perforget}
\end{table}

\begin{figure}[t]
\centering
\includegraphics[width=0.95\columnwidth]{fig_per_forget_profile.pdf}
\caption{\textbf{Per-forget-set variance ($n{=}100$).} Box + jitter of geo-mean $r$ for each layer; one dot per forget set. Spread $\max/\min$ reaches $1.59\times$ on $\Lone$, $1.50\times$ on $\Ltwo$, $1.13\times$ on $\Lthree$. \texttt{triplet\_073} (``storm'') is the worst on $\Lone$ and $\Ltwo$ but not on $\Lthree$ (where \texttt{triplet\_072} is the maximum, $1.26\times$); \texttt{triplet\_029} is the $\Lone$ minimum. Distribution shape is markedly tighter as the layer recedes from $\Df$.}
\label{fig:per-forget-profile}
\end{figure}

\textbf{(C3) Side-effect is a property of $\Df$.} Under identical optimisation, switching $\Df$ changes $\Lone$ by $1.59\times$ end-to-end (log-scale spread $\approx 1.7$). On a one-page reading: the question \emph{``how much will unlearning hurt?''} has no single answer for a given algorithm; it depends on which $\Df$ is unlearned. This is the empirical wedge that motivates the audit problem in Sec.~\ref{sec:audit}.
